\documentclass[12pt]{article}
\usepackage[T1]{fontenc}
\usepackage[utf8]{inputenc}
\usepackage{lmodern}
\usepackage{textcomp}
\usepackage{enumitem}
\usepackage{geometry}
\usepackage{graphicx}
\usepackage{amsmath, amssymb}
\geometry{margin=1in}
\begin{document}
\begin{center}
Constitutional Law - Undergraduate Level Assessment 3 \
Total Marks: 40 \
Answer all questions.
\end{center}

\section*{Multiple Choice (10 marks)}
\begin{enumerate}[label=\arabic*.]
\item The \_\_\_\_\_\_\_\_ School of jurisprudence maintains that the law is shaped by logic.
\begin{enumerate}[label=(\alph*)]
    \item Historical
    \item Analytical
    \item Command
    \item Sociological
\end{enumerate}
\item What is the purpose of politics, according to Aristotle?
\begin{enumerate}[label=(\alph*)]
    \item To advance the interests of politicians.
    \item To produce virtuous citizens and encourage righteousness in individuals.
    \item To obtain power.
    \item To secure justice by abolishing slavery.
\end{enumerate}
\item Why is there an increasing recognition that retribution is a fundamental element of punishment?
\begin{enumerate}[label=(\alph*)]
    \item Because rehabilitation of offenders presumes a recognition of the role of vengeance.
    \item Because the lex talionis is misguided.
    \item Because other justifications appear to have failed.
    \item Because revenge is anachronistic.
\end{enumerate}
\item Hart's analysis of law distinguishes between
\begin{enumerate}[label=(\alph*)]
    \item Cause and effect
    \item Theory and fact
    \item Being obliged and having an obligation
    \item Corporeal and incorporeal rights
\end{enumerate}
\item Which rule of jus cogens was the first to be accepted explicitly as such by the ICJ?
\begin{enumerate}[label=(\alph*)]
    \item The prohibition of the use of force
    \item The prohibition of torture
    \item The prohibition of genocide
    \item The principle of self-determination
\end{enumerate}
\end{enumerate}

\section*{True / False (10 marks)}
\textit{Answer True or False}
\begin{enumerate}[label=\arabic*.]
\setcounter{enumi}{5}
\item True or False: John Locke dismissed the idea of natural rights as irrelevant to governance.
\item True or False: Recognition is simply declaratory of statehood and does not determine the existence of a state.
\item True or False: Principles describe rights, while policies describe goals within Dworkin's framework of legal interpretation.
\item True or False: Leopold Pospisil identifies 'dispute' as one of the four essential elements manifested by law.
\item True or False: Olivecrona's argument is well-supported by extensive empirical evidence demonstrating the psychological basis of law.
\end{enumerate}

\section*{Long Form Response (20 marks)}
\textit{Answer each question with a clear and complete explanation.}
\begin{enumerate}[label=\arabic*.]
\setcounter{enumi}{10}
\item Discuss Rousseau's assertion that the "Social Contract is not a historical fact but a hypothetical construction of reason." Analyze its implications for constitutional law and modern governance.
\item Discuss the implications of allowing a reservation to the definition of torture in the ICCPR, analyzing whether it aligns with the treaty's object and purpose in contemporary legal practice.
\end{enumerate}

\vfill
\noindent\textit{End of Paper}
\end{document}